\documentclass[10pt,a4paper]{amsart}
\usepackage{amsmath,amssymb,amsthm,graphicx}
% For commutative diagrams we use Xy-pic
\usepackage[all]{xy}

% We use 2cell for 2-commutative diagrams.
\xyoption{2cell}
% Theorem environments.
%
\theoremstyle{plain}
\newtheorem{theorem}[subsection]{Theorem}
\newtheorem{proposition}[subsection]{Proposition}
\newtheorem{lemma}[subsection]{Lemma}

\theoremstyle{definition}
\newtheorem{definition}[subsection]{Definition}
\newtheorem{example}[subsection]{Example}
\newtheorem{exercise}[subsection]{Exercise}
\newtheorem{situation}[subsection]{Situation}

\theoremstyle{remark}
\newtheorem{remark}[subsection]{Remark}
\newtheorem{remarks}[subsection]{Remarks}

\numberwithin{equation}{subsection}

% Macros
%
\def\lim{\mathop{\mathrm{lim}}\nolimits}
\def\colim{\mathop{\mathrm{colim}}\nolimits}
\def\Spec{\mathop{\mathrm{Spec}}}
\def\Hom{\mathop{\mathrm{Hom}}\nolimits}
\def\Ext{\mathop{\mathrm{Ext}}\nolimits}
\def\SheafHom{\mathop{\mathcal{H}\!\mathit{om}}\nolimits}
\def\SheafExt{\mathop{\mathcal{E}\!\mathit{xt}}\nolimits}
\def\Sch{\mathit{Sch}}
\def\Mor{\mathop{\mathrm{Mor}}\nolimits}
\def\Ob{\mathop{\mathrm{Ob}}\nolimits}
\def\Sh{\mathop{\mathit{Sh}}\nolimits}
\def\NL{\mathop{N\!L}\nolimits}
\def\CH{\mathop{\mathrm{CH}}\nolimits}
\def\proetale{{pro\text{-}\acute{e}tale}}
\def\etale{{\acute{e}tale}}
\def\QCoh{\mathit{QCoh}}
\def\Ker{\mathop{\mathrm{Ker}}}
\def\Im{\mathop{\mathrm{Im}}}
\def\Coker{\mathop{\mathrm{Coker}}}
\def\Coim{\mathop{\mathrm{Coim}}}

% Boxtimes
%
\DeclareMathSymbol{\boxtimes}{\mathbin}{AMSa}{"02}

%
% Macros for moduli stacks/spaces
%
\def\QCohstack{\mathcal{QC}\!\mathit{oh}}
\def\Cohstack{\mathcal{C}\!\mathit{oh}}
\def\Spacesstack{\mathcal{S}\!\mathit{paces}}
\def\Quotfunctor{\mathrm{Quot}}
\def\Hilbfunctor{\mathrm{Hilb}}
\def\Curvesstack{\mathcal{C}\!\mathit{urves}}
\def\Polarizedstack{\mathcal{P}\!\mathit{olarized}}
\def\Complexesstack{\mathcal{C}\!\mathit{omplexes}}
% \Pic is the operator that assigns to X its picard group, usage \Pic(X)
% \Picardstack_{X/B} denotes the Picard stack of X over B
% \Picardfunctor_{X/B} denotes the Picard functor of X over B
\def\Pic{\mathop{\mathrm{Pic}}\nolimits}
\def\Picardstack{\mathcal{P}\!\mathit{ic}}
\def\Picardfunctor{\mathrm{Pic}}
\def\Deformationcategory{\mathcal{D}\!\mathit{ef}}
\newenvironment{reference}{}{}
\newenvironment{slogan}{}{}
\newenvironment{history}{}{}
\usepackage{cleveref}
\crefname{originalproposition}{original proposition}{original propositions}
\Crefname{originalproposition}{Original proposition}{Original propositions}
\makeatletter
\newlabel{proposition-enough-weakly-contractibles}{{099Q}{source}{}{sourcepointer-1}{}}
\newlabel{proposition-enough-weakly-contractibles@cref}{{[originalproposition][0][]099Q}{[1][0][]source}{}{}{}}
\makeatother
\crefname{originallemma}{original lemma}{original lemmas}
\Crefname{originallemma}{Original lemma}{Original lemmas}
\makeatletter
\newlabel{lemma-naive-completion}{{099M}{source}{}{sourcepointer-2}{}}
\newlabel{lemma-naive-completion@cref}{{[originallemma][0][]099M}{[1][0][]source}{}{}{}}
\makeatother
\crefname{originaldefinition}{original definition}{original definitions}
\Crefname{originaldefinition}{Original definition}{Original definitions}
\makeatletter
\newlabel{etale-cohomology-definition-ctf}{{03TQ}{source}{}{sourcepointer-3}{}}
\newlabel{etale-cohomology-definition-ctf@cref}{{[originaldefinition][0][]03TQ}{[1][0][]source}{}{}{}}
\makeatother
\crefname{originallemma}{original lemma}{original lemmas}
\Crefname{originallemma}{Original lemma}{Original lemmas}
\makeatletter
\newlabel{sites-cohomology-lemma-bounded}{{0942}{source}{}{sourcepointer-4}{}}
\newlabel{sites-cohomology-lemma-bounded@cref}{{[originallemma][0][]0942}{[1][0][]source}{}{}{}}
\makeatother
\crefname{originalproposition}{original proposition}{original propositions}
\Crefname{originalproposition}{Original proposition}{Original propositions}
\makeatletter
\newlabel{sites-cohomology-proposition-enough-weakly-contractibles}{{0947}{source}{}{sourcepointer-5}{}}
\newlabel{sites-cohomology-proposition-enough-weakly-contractibles@cref}{{[originalproposition][0][]0947}{[1][0][]source}{}{}{}}
\makeatother
\crefname{originallemma}{original lemma}{original lemmas}
\Crefname{originallemma}{Original lemma}{Original lemmas}
\makeatletter
\newlabel{lemma-compare-truncations-constructible}{{09BR}{source}{}{sourcepointer-6}{}}
\newlabel{lemma-compare-truncations-constructible@cref}{{[originallemma][0][]09BR}{[1][0][]source}{}{}{}}
\makeatother
\crefname{originallemma}{original lemma}{original lemmas}
\Crefname{originallemma}{Original lemma}{Original lemmas}
\makeatletter
\newlabel{lemma-compare-constructible-derived}{{09AL}{source}{}{sourcepointer-7}{}}
\newlabel{lemma-compare-constructible-derived@cref}{{[originallemma][0][]09AL}{[1][0][]source}{}{}{}}
\makeatother
\crefname{originallemma}{original lemma}{original lemmas}
\Crefname{originallemma}{Original lemma}{Original lemmas}
\makeatletter
\newlabel{sites-modules-lemma-morphism-locally-constant}{{093S}{source}{}{sourcepointer-8}{}}
\newlabel{sites-modules-lemma-morphism-locally-constant@cref}{{[originallemma][0][]093S}{[1][0][]source}{}{}{}}
\makeatother
\crefname{originallemma}{original lemma}{original lemmas}
\Crefname{originallemma}{Original lemma}{Original lemmas}
\makeatletter
\newlabel{algebra-lemma-finite-projective}{{00NX}{source}{}{sourcepointer-9}{}}
\newlabel{algebra-lemma-finite-projective@cref}{{[originallemma][0][]00NX}{[1][0][]source}{}{}{}}
\makeatother
\crefname{originallemma}{original lemma}{original lemmas}
\Crefname{originallemma}{Original lemma}{Original lemmas}
\makeatletter
\newlabel{algebra-lemma-lift-idempotents-noncommutative}{{05BU}{source}{}{sourcepointer-10}{}}
\newlabel{algebra-lemma-lift-idempotents-noncommutative@cref}{{[originallemma][0][]05BU}{[1][0][]source}{}{}{}}
\makeatother
\crefname{originallemma}{original lemma}{original lemmas}
\Crefname{originallemma}{Original lemma}{Original lemmas}
\makeatletter
\newlabel{sites-modules-lemma-completion-flat}{{093M}{source}{}{sourcepointer-11}{}}
\newlabel{sites-modules-lemma-completion-flat@cref}{{[originallemma][0][]093M}{[1][0][]source}{}{}{}}
\makeatother
\crefname{originalsection}{original section}{original sections}
\Crefname{originalsection}{Original section}{Original sections}
\makeatletter
\newlabel{section-derived-completion-noetherian}{{099L}{source}{}{sourcepointer-12}{}}
\newlabel{section-derived-completion-noetherian@cref}{{[originalsection][0][]099L}{[1][0][]source}{}{}{}}
\makeatother
\crefname{originaldefinition}{original definition}{original definitions}
\Crefname{originaldefinition}{Original definition}{Original definitions}
\makeatletter
\newlabel{sites-cohomology-definition-tor-amplitude}{{08FZ}{source}{}{sourcepointer-13}{}}
\newlabel{sites-cohomology-definition-tor-amplitude@cref}{{[originaldefinition][0][]08FZ}{[1][0][]source}{}{}{}}
\makeatother
\title{Finite projective local models for constructible complete complexes on weakly contractible affine schemes}
\author{The Stacks Project Authors}
\date{stacks-project:proetale.tex:local-snapshot-20261011; GNU FDL 1.2 or later}
\begin{document}
\maketitle
\noindent\textit{Editorial note.} D\_cons is defined in the included native context by derived completeness and constructible finiteTor-dimension reduction. The theorem additionally assumes locally constant reduced cohomology and weakly contractible affineX.\par\medskip
\noindent\textit{Allowed standard prerequisites (editorial).}\begin{itemize}
\item Explicit prerequisite boundary. Pro-etale sites, weakly contractible affine schemes, derivedI-completion, constructible complexes with finiteTor amplitude, lifting idempotents and finite projective modules. The local complete-zero and complete-limit support is retained.
\end{itemize}\medskip
\ifdefined\hypertarget\hypertarget{sourcepointer-1}{}\fi
\noindent\textit{Source reference.} Original proposition 099Q: focused-seeds/AGP/sources/stacks-snapshot-20261011/tags.txt, source lines 11733--11733. Stable Stacks tag for an expressly allowed foundational or notational input; the counted construction remains the two concrete core obligations. The retained local core is not replaced by this pointer. See standard\_prerequisites.\par\medskip
\ifdefined\hypertarget\hypertarget{sourcepointer-2}{}\fi
\noindent\textit{Source reference.} Original lemma 099M: focused-seeds/AGP/sources/stacks-snapshot-20261011/tags.txt, source lines 11730--11730. Stable Stacks tag for an expressly allowed foundational or notational input; the counted construction remains the two concrete core obligations. The retained local core is not replaced by this pointer. See standard\_prerequisites.\par\medskip
\ifdefined\hypertarget\hypertarget{sourcepointer-3}{}\fi
\noindent\textit{Source reference.} Original definition 03TQ: focused-seeds/AGP/sources/stacks-snapshot-20261011/tags.txt, source lines 5013--5013. Stable Stacks tag for an expressly allowed foundational or notational input; the counted construction remains the two concrete core obligations. The retained local core is not replaced by this pointer. See standard\_prerequisites.\par\medskip
\ifdefined\hypertarget\hypertarget{sourcepointer-4}{}\fi
\noindent\textit{Source reference.} Original lemma 0942: focused-seeds/AGP/sources/stacks-snapshot-20261011/tags.txt, source lines 11535--11535. Stable Stacks tag for an expressly allowed foundational or notational input; the counted construction remains the two concrete core obligations. The retained local core is not replaced by this pointer. See standard\_prerequisites.\par\medskip
\ifdefined\hypertarget\hypertarget{sourcepointer-5}{}\fi
\noindent\textit{Source reference.} Original proposition 0947: focused-seeds/AGP/sources/stacks-snapshot-20261011/tags.txt, source lines 11540--11540. Stable Stacks tag for an expressly allowed foundational or notational input; the counted construction remains the two concrete core obligations. The retained local core is not replaced by this pointer. See standard\_prerequisites.\par\medskip
\ifdefined\hypertarget\hypertarget{sourcepointer-6}{}\fi
\noindent\textit{Source reference.} Original lemma 09BR: focused-seeds/AGP/sources/stacks-snapshot-20261011/tags.txt, source lines 11804--11804. Stable Stacks tag for an expressly allowed foundational or notational input; the counted construction remains the two concrete core obligations. The retained local core is not replaced by this pointer. See standard\_prerequisites.\par\medskip
\ifdefined\hypertarget\hypertarget{sourcepointer-7}{}\fi
\noindent\textit{Source reference.} Original lemma 09AL: focused-seeds/AGP/sources/stacks-snapshot-20261011/tags.txt, source lines 11764--11764. Stable Stacks tag for an expressly allowed foundational or notational input; the counted construction remains the two concrete core obligations. The retained local core is not replaced by this pointer. See standard\_prerequisites.\par\medskip
\ifdefined\hypertarget\hypertarget{sourcepointer-8}{}\fi
\noindent\textit{Source reference.} Original lemma 093S: focused-seeds/AGP/sources/stacks-snapshot-20261011/tags.txt, source lines 11525--11525. Stable Stacks tag for an expressly allowed foundational or notational input; the counted construction remains the two concrete core obligations. The retained local core is not replaced by this pointer. See standard\_prerequisites.\par\medskip
\ifdefined\hypertarget\hypertarget{sourcepointer-9}{}\fi
\noindent\textit{Source reference.} Original lemma 00NX: focused-seeds/AGP/sources/stacks-snapshot-20261011/tags.txt, source lines 875--875. Stable Stacks tag for an expressly allowed foundational or notational input; the counted construction remains the two concrete core obligations. The retained local core is not replaced by this pointer. See standard\_prerequisites.\par\medskip
\ifdefined\hypertarget\hypertarget{sourcepointer-10}{}\fi
\noindent\textit{Source reference.} Original lemma 05BU: focused-seeds/AGP/sources/stacks-snapshot-20261011/tags.txt, source lines 6897--6897. Stable Stacks tag for an expressly allowed foundational or notational input; the counted construction remains the two concrete core obligations. The retained local core is not replaced by this pointer. See standard\_prerequisites.\par\medskip
\ifdefined\hypertarget\hypertarget{sourcepointer-11}{}\fi
\noindent\textit{Source reference.} Original lemma 093M: focused-seeds/AGP/sources/stacks-snapshot-20261011/tags.txt, source lines 11520--11520. Stable Stacks tag for an expressly allowed foundational or notational input; the counted construction remains the two concrete core obligations. The retained local core is not replaced by this pointer. See standard\_prerequisites.\par\medskip
\ifdefined\hypertarget\hypertarget{sourcepointer-12}{}\fi
\noindent\textit{Source reference.} Original section 099L: focused-seeds/AGP/sources/stacks-snapshot-20261011/tags.txt, source lines 11729--11729. Stable Stacks tag for an expressly allowed foundational or notational input; the counted construction remains the two concrete core obligations. The retained local core is not replaced by this pointer. See standard\_prerequisites.\par\medskip
\ifdefined\hypertarget\hypertarget{sourcepointer-13}{}\fi
\noindent\textit{Source reference.} Original definition 08FZ: focused-seeds/AGP/sources/stacks-snapshot-20261011/tags.txt, source lines 10727--10727. Stable Stacks tag for an expressly allowed foundational or notational input; the counted construction remains the two concrete core obligations. The retained local core is not replaced by this pointer. See standard\_prerequisites.\par\medskip
\begin{lemma}
\label{lemma-local-structure-constructible-complex}
Let $X$ be a weakly contractible affine scheme. Let $\Lambda$ be a Noetherian
ring and let $I \subset \Lambda$ be an ideal. Let $K$ be an object of
$D_{cons}(X, \Lambda)$ such that the cohomology sheaves of
$K \otimes_\Lambda^\mathbf{L} \underline{\Lambda/I}$ are locally
constant. Then there exists a finite disjoint open covering
$X = \coprod U_i$ and for each $i$ a finite collection of
finite projective $\Lambda^\wedge$-modules $M^a, \ldots, M^b$
such that $K|_{U_i}$ is represented by a complex
$$
(\underline{M^a})^\wedge \to \ldots \to (\underline{M^b})^\wedge
$$
in $D(U_{i, \proetale}, \Lambda)$ for some maps of sheaves of
$\Lambda$-modules $(\underline{M^i})^\wedge \to (\underline{M^{i + 1}})^\wedge$.
\end{lemma}

\section{A suitable derived category}
\label{section-suitable-derived}

\noindent
Let $X$ be a scheme. It will turn out that for many schemes $X$ a suitable
derived category of $\ell$-adic sheaves can be gotten by considering the
derived complete objects $K$ of $D(X_\proetale, \Lambda)$ with the property
that $K \otimes_\Lambda^\mathbf{L} \mathbf{F}_\ell$ is bounded with
constructible cohomology sheaves. Here is the general definition.

\begin{definition}
\label{definition-Dbc}
Let $\Lambda$ be a Noetherian ring and let $I \subset \Lambda$ be an ideal.
Let $X$ be a scheme. An object $K$ of $D(X_\proetale, \Lambda)$ is called
{\it constructible} if
\begin{enumerate}
\item $K$ is derived complete with respect to $I$,
\item $K \otimes_\Lambda^\mathbf{L} \underline{\Lambda/I}$
has constructible cohomology sheaves and locally has finite tor dimension.
\end{enumerate}
We denote $D_{cons}(X, \Lambda)$ the full subcategory of constructible
$K$ in $D(X_\proetale, \Lambda)$.
\end{definition}

\noindent
Recall that with our conventions a complex of finite tor dimension
is bounded (Cohomology on Sites, Definition
\ref{sites-cohomology-definition-tor-amplitude}).
In fact, let's collect everything proved so far in a lemma.

\begin{lemma}
\label{lemma-derived-complete-zero}
Let $X$ be a scheme. Let $\Lambda$ be a ring and let
$I \subset \Lambda$ be a finitely generated ideal.
Let $\mathcal{F}$ be a sheaf of $\Lambda$-modules on $X_\proetale$.
If $\mathcal{F}$ is derived complete and $\mathcal{F}/I\mathcal{F} = 0$,
then $\mathcal{F} = 0$.
\end{lemma}

\begin{proof}
Assume that $\mathcal{F}/I\mathcal{F}$ is zero.
Let $I = (f_1, \ldots, f_r)$. Let $i < r$ be the largest
integer such that $\mathcal{G} = \mathcal{F}/(f_1, \ldots, f_i)\mathcal{F}$
is nonzero. If $i$ does not exist, then $\mathcal{F} = 0$ which is what we
want to show. Then $\mathcal{G}$ is derived complete as a cokernel
of a map between derived complete modules, see
Proposition \ref{proposition-enough-weakly-contractibles}.
By our choice of $i$ we have that $f_{i + 1} : \mathcal{G} \to \mathcal{G}$
is surjective. Hence
$$
\lim (\ldots \to \mathcal{G} \xrightarrow{f_{i + 1}} \mathcal{G}
\xrightarrow{f_{i + 1}} \mathcal{G})
$$
is nonzero, contradicting the derived completeness of $\mathcal{G}$.
\end{proof}

\begin{lemma}
\label{lemma-derived-complete-limit}
Let $X$ be a weakly contractible affine scheme.
Let $\Lambda$ be a Noetherian ring and let $I \subset \Lambda$ be an ideal.
Let $\mathcal{F}$ be a derived complete sheaf of $\Lambda$-modules
on $X_\proetale$ with $\mathcal{F}/I\mathcal{F}$ a locally
constant sheaf of $\Lambda/I$-modules of finite type.
Then there exists an integer $t$ and a surjective map
$$
(\underline{\Lambda}^\wedge)^{\oplus t} \to \mathcal{F}
$$
\end{lemma}

\begin{proof}
Since $X$ is weakly contractible, there exists a finite disjoint open
covering $X = \coprod U_i$ such that $\mathcal{F}/I\mathcal{F}|_{U_i}$
is isomorphic to the constant sheaf associated to a finite $\Lambda/I$-module
$M_i$. Choose finitely many generators $m_{ij}$ of $M_i$. We
can find sections $s_{ij} \in \mathcal{F}(X)$ restricting to
$m_{ij}$ viewed as a section of $\mathcal{F}/I\mathcal{F}$ over $U_i$. 
Let $t$ be the total number of $s_{ij}$. Then we obtain a map
$$
\alpha : \underline{\Lambda}^{\oplus t} \longrightarrow \mathcal{F}
$$
which is surjective modulo $I$ by construction. By
Lemma \ref{lemma-naive-completion}
the derived completion of $\underline{\Lambda}^{\oplus t}$ is the
sheaf $(\underline{\Lambda}^\wedge)^{\oplus t}$. Since $\mathcal{F}$
is derived complete we see that $\alpha$ factors through a map
$$
\alpha^\wedge :
(\underline{\Lambda}^\wedge)^{\oplus t}
\longrightarrow
\mathcal{F}
$$
Then $\mathcal{Q} = \Coker(\alpha^\wedge)$ is a derived complete
sheaf of $\Lambda$-modules by
Proposition \ref{proposition-enough-weakly-contractibles}.
By construction $\mathcal{Q}/I\mathcal{Q} = 0$. It follows from
Lemma \ref{lemma-derived-complete-zero}
that $\mathcal{Q} = 0$ which is what we wanted to show.
\end{proof}

\begin{lemma}
\label{lemma-describe-constructible-complexes}
In the situation above suppose $K$ is in $D_{cons}(X, \Lambda)$
and $X$ is quasi-compact. Set
$K_n = K \otimes_\Lambda^\mathbf{L} \underline{\Lambda/I^n}$.
There exist $a, b$ such that
\begin{enumerate}
\item $K = R\lim K_n$ and $H^i(K) = 0$ for $i \not \in [a, b]$,
\item each $K_n$ has tor amplitude in $[a, b]$,
\item each $K_n$ has constructible cohomology sheaves,
\item each $K_n = \epsilon^{-1}L_n$ for some
$L_n \in D_{ctf}(X_\etale, \Lambda/I^n)$
(\'Etale Cohomology, Definition \ref{etale-cohomology-definition-ctf}).
\end{enumerate}
\end{lemma}

\begin{proof}
Since $K$ is derived complete, we have $K = R\lim K_n$ by
Lemma \ref{lemma-naive-completion}.
By definition of locally having finite tor dimension, we can find
$a, b$ such that $K_1$ has tor amplitude in $[a, b]$.
Part (2) follows from 
Cohomology on Sites, Lemma \ref{sites-cohomology-lemma-bounded}.
Then (1) follows as $K$ is derived complete by the description
of limits in
Cohomology on Sites, Proposition
\ref{sites-cohomology-proposition-enough-weakly-contractibles}
and the fact that $H^b(K_{n + 1}) \to H^b(K_n)$ is surjective
as $K_n = K_{n + 1} \otimes^\mathbf{L}_\Lambda \underline{\Lambda/I^n}$.
Part (3) follows from
Lemma \ref{lemma-compare-truncations-constructible},
Part (4) follows from
Lemma \ref{lemma-compare-constructible-derived}
and the fact that $L_n$ has finite tor dimension because $K_n$ does
(small argument omitted).
\end{proof}

\begin{proof}
We freely use the results of
Lemma \ref{lemma-describe-constructible-complexes}.
Choose $a, b$ as in that lemma. We will prove the lemma by
induction on $b - a$. Let $\mathcal{F} = H^b(K)$.
Note that $\mathcal{F}$ is a derived complete sheaf of
$\Lambda$-modules by
Proposition \ref{proposition-enough-weakly-contractibles}.
Moreover $\mathcal{F}/I\mathcal{F}$ is a locally
constant sheaf of $\Lambda/I$-modules of finite type.
Apply Lemma \ref{lemma-derived-complete-limit} to get a surjection
$\rho : (\underline{\Lambda}^\wedge)^{\oplus t} \to \mathcal{F}$.

\medskip\noindent
If $a = b$, then $K = \mathcal{F}[-b]$. In this case we see that
$$
\mathcal{F} \otimes_\Lambda^\mathbf{L} \underline{\Lambda/I} =
\mathcal{F}/I\mathcal{F}
$$
As $X$ is weakly contractible and $\mathcal{F}/I\mathcal{F}$
locally constant, we can find a finite disjoint
union decomposition $X = \coprod U_i$ by affine opens $U_i$
and $\Lambda/I$-modules $\overline{M}_i$
such that $\mathcal{F}/I\mathcal{F}$ restricts to
$\underline{\overline{M}_i}$ on $U_i$. After refining the covering
we may assume the map
$$
\rho|_{U_i} \bmod I :
\underline{\Lambda/I}^{\oplus t}
\longrightarrow
\underline{\overline{M}_i}
$$
is equal to $\underline{\alpha_i}$ for some surjective module map
$\alpha_i : \Lambda/I^{\oplus t} \to \overline{M}_i$, see
Modules on Sites, Lemma \ref{sites-modules-lemma-morphism-locally-constant}.
Note that each $\overline{M}_i$ is a finite $\Lambda/I$-module.
Since $\mathcal{F}/I\mathcal{F}$ has tor amplitude in $[0, 0]$
we conclude that $\overline{M}_i$ is a flat $\Lambda/I$-module.
Hence $\overline{M}_i$ is finite projective
(Algebra, Lemma \ref{algebra-lemma-finite-projective}).
Hence we can find a projector
$\overline{p}_i : (\Lambda/I)^{\oplus t} \to (\Lambda/I)^{\oplus t}$
whose image maps isomorphically to $\overline{M}_i$ under the map $\alpha_i$.
We can lift $\overline{p}_i$ to a projector
$p_i : (\Lambda^\wedge)^{\oplus t} \to
(\Lambda^\wedge)^{\oplus t}$\footnote{Proof: by
Algebra, Lemma \ref{algebra-lemma-lift-idempotents-noncommutative}
we can lift $\overline{p}_i$ to a compatible system of
projectors $p_{i, n} : (\Lambda/I^n)^{\oplus t} \to (\Lambda/I^n)^{\oplus t}$
and then we set $p_i = \lim p_{i, n}$ which works because
$\Lambda^\wedge = \lim \Lambda/I^n$.}.
Then $M_i = \Im(p_i)$ is a finite $I$-adically complete
$\Lambda^\wedge$-module with $M_i/IM_i = \overline{M}_i$.
Over $U_i$ consider the maps
$$
\underline{M_i}^\wedge \to
(\underline{\Lambda}^\wedge)^{\oplus t} \to
\mathcal{F}|_{U_i}
$$
By construction the composition induces an isomorphism modulo $I$.
The source and target are derived complete, hence so are the cokernel
$\mathcal{Q}$ and the kernel $\mathcal{K}$. We have
$\mathcal{Q}/I\mathcal{Q} = 0$ by construction hence $\mathcal{Q}$
is zero by Lemma \ref{lemma-derived-complete-zero}.
Then
$$
0 \to \mathcal{K}/I\mathcal{K} \to
\underline{\overline{M}_i}
\to \mathcal{F}/I\mathcal{F} \to 0
$$
is exact by the vanishing of $\text{Tor}_1$ see at the start of this
paragraph; also use that
$\underline{\Lambda}^\wedge/I\overline{\Lambda}^\wedge$ by
Modules on Sites, Lemma \ref{sites-modules-lemma-completion-flat}
to see that
$\underline{M_i}^\wedge/I\underline{M_i}^\wedge = \underline{\overline{M}_i}$.
Hence $\mathcal{K}/I\mathcal{K} = 0$ by construction and we conclude
that $\mathcal{K} = 0$ as before. This proves the result in case $a = b$.

\medskip\noindent
If $b > a$, then we lift the map $\rho$ to a map
$$
\tilde \rho : (\underline{\Lambda}^\wedge)^{\oplus t}[-b] \longrightarrow K
$$
in $D(X_\proetale, \Lambda)$. This is possible as we can think of
$K$ as a complex of $\underline{\Lambda}^\wedge$-modules by
discussion in the introduction to
Section \ref{section-derived-completion-noetherian}
and because $X_\proetale$ is weakly contractible
hence there is no obstruction to lifting the elements
$\rho(e_s) \in H^0(X, \mathcal{F})$ to elements of $H^b(X, K)$.
Fitting $\tilde \rho$ into a distinguished triangle
$$
(\underline{\Lambda}^\wedge)^{\oplus t}[-b] \to K \to L \to
(\underline{\Lambda}^\wedge)^{\oplus t}[-b + 1]
$$
we see that $L$ is an object of $D_{cons}(X, \Lambda)$ such
that $L \otimes_\Lambda^\mathbf{L} \underline{\Lambda/I}$
has tor amplitude contained in $[a, b - 1]$ (details omitted).
By induction we can describe $L$ locally as stated in the lemma, say
$L$ is isomorphic to
$$
(\underline{M^a})^\wedge \to \ldots \to (\underline{M^{b - 1}})^\wedge
$$
The map
$L \to (\underline{\Lambda}^\wedge)^{\oplus t}[-b + 1]$
corresponds to a map
$(\underline{M^{b - 1}})^\wedge \to (\underline{\Lambda}^\wedge)^{\oplus t}$
which allows us to extend the complex by one. The corresponding
complex is isomorphic to $K$ in the derived category by the properties
of triangulated categories. This finishes the proof.
\end{proof}

\bibliographystyle{plain}
\bibliography{my}
\end{document}
